\documentclass[11pt]{article}
\usepackage[T1]{fontenc}
\usepackage{textcomp}
\usepackage[margin=0.75in]{geometry}
\usepackage{amsmath,amssymb,amsthm}
\usepackage{array,booktabs}
\usepackage{hyperref}
\setlength{\parindent}{0pt}
\newtheorem{theorem}{Theorem}
\theoremstyle{definition}
\newtheorem{definition}{Definition}
\title{Flow Proof Artifact: prop-13-angles-on-straight-line}
\author{Generated by \texttt{flow doc proof}}
\date{Flow Proof Book}
\begin{document}
\maketitle
If a straight line stands on a straight line, the adjacent angles equal two right angles.\\[0.5em]
\textbf{Source.} euclid — Elements, Book I, Proposition 13\\[0.5em]
\bigskip
\noindent{\large \textbf{Axiom 1.} \textit{If a straight line stands on a straight line, the adjacent angles equal two right angles} \hfill the Euclidean plane \cdot Euclid Book I \cdot Proposition 13: adjacent angles on a straight line sum to two right angles}\par
\smallskip\noindent\textit{Coordinate: the Euclidean plane \cdot Euclid Book I \cdot Proposition 13: adjacent angles on a straight line sum to two right angles}\par
\smallskip\noindent\textit{Source: euclid — Elements, Book I, Proposition 13}\par
\medskip
\noindent\textbf{Goal.} If a straight line stands on a straight line, the adjacent angles equal two right angles\par
\noindent$\displaystyle \angle APC + \angle APB = 180^\circ$\par
\medskip
\noindent
\begin{tabular}{@{} >{\raggedright\arraybackslash}p{0.06\textwidth} >{\raggedright\arraybackslash}p{0.47\textwidth} >{\raggedleft\arraybackslash}p{0.41\textwidth} @{}}
\textbf{\#} & \textbf{Proof} & \textbf{Mathematics} \\
\midrule
\textbf{1.} & We accept proposition 13: adjacent angles on a straight line sum to two right angles for Book I of the Elements in the Euclidean plane without proof — an ontological commitment, not a lemma. & \\[0.45em]
\textbf{2.} & We invoke the derived fact: line\_PC\_stands\_on\_line\_AB\_at\_P. & \\[0.45em]
\textbf{3.} & We invoke the derived fact: Common Notion 4: coinciding straight lines have equal adjacent angles. & \\[0.45em]
\textbf{4.} & From step 2 and step 3, by the chain of results established in Book I (step 2 and step 3), we obtain angle APC plus angle APB equals two right angles. Hence proven. & $\displaystyle \angle APC + \angle APB = 180^\circ$ \\[0.45em]
\bottomrule
\end{tabular}
\label{thm:Geometry-Euclid Book I-proposition_13_adjacent_angles_on_a_straight_line_sum_to_two_right_angles}
\par\medskip\hrule
\end{document}
